\documentclass[11pt]{article}
% ===========================================================================
%  Shared preamble for all MPM2D worksheets — input right after \documentclass.
%  (Files beginning with "_" are skipped by mpm2d-worksheet-publish.mjs.)
% ===========================================================================
\usepackage[letterpaper,top=2.2cm,bottom=1.9cm,left=1.6cm,right=1.6cm,headheight=28pt,footskip=24pt]{geometry}
\usepackage{amsmath,amssymb}
\usepackage{cancel}
\usepackage[T1]{fontenc}
\usepackage[utf8]{inputenc}
\usepackage{lmodern}
\usepackage{xcolor}
\usepackage[most]{tcolorbox}
\usepackage{enumitem}
\usepackage{multicol}
\usepackage{fancyhdr}
\usepackage{tikz}
\usetikzlibrary{calc}
\usepackage{pgfplots}
\pgfplotsset{compat=1.18}

\definecolor{exblue}{HTML}{4A90E2}
\definecolor{exbg}{HTML}{E6F3FF}
\definecolor{qorange}{HTML}{E69138}
\definecolor{qbg}{HTML}{FFF7CC}
\definecolor{keygray}{HTML}{475569}
\definecolor{keybg}{HTML}{F1F5F9}
\definecolor{iagreen}{HTML}{14653B}
\definecolor{titleblue}{HTML}{1D4ED8}
% Rotating palette so each worked example gets its own colour.
\definecolor{ex0}{HTML}{2563A0}\definecolor{ex1}{HTML}{3B7D3B}\definecolor{ex2}{HTML}{8A5A00}
\definecolor{ex3}{HTML}{A3327A}\definecolor{ex4}{HTML}{6D28A3}\definecolor{ex5}{HTML}{B08900}
\definecolor{ex6}{HTML}{0E7490}\definecolor{ex7}{HTML}{9A3412}\definecolor{ex8}{HTML}{15803D}

\pagestyle{fancy}
\fancyhf{}
\fancyhead[L]{\footnotesize NAME: \underline{\hspace{3.4cm}}}
\fancyhead[C]{\textbf{Integration Academy}\\[-2pt]{\scriptsize Math Simplified \textperiodcentered{} Success Amplified}}
\fancyhead[R]{\footnotesize MPM2D}
\fancyfoot[L]{\scriptsize dr.merie@IntegrationAcademy.ca}
\fancyfoot[C]{\scriptsize\thepage}
\fancyfoot[R]{\scriptsize IntegrationAcademy.ca}
\renewcommand{\headrulewidth}{1.2pt}
\renewcommand{\footrulewidth}{0.4pt}
\setlength{\parindent}{0pt}
\setlength{\parskip}{4pt}

% Examples are NOT breakable, so each example + its solution stays on one page.
\newtcolorbox{example}[2]{colframe=#1,colback=#1!7,coltitle=white,colbacktitle=#1,fonttitle=\bfseries,title={#2},arc=3pt,boxrule=1pt}
\newtcolorbox{qbox}[1]{colframe=qorange,colback=qbg,coltitle=white,colbacktitle=qorange,fonttitle=\bfseries,title={#1},arc=3pt,breakable,boxrule=1pt}
\newtcolorbox{keybox}{colframe=keygray,colback=keybg,coltitle=white,colbacktitle=keygray,fonttitle=\bfseries,title={$\checkmark$\ Answer Key},arc=3pt,breakable,boxrule=1pt}
\newtcolorbox{recapbox}{colframe=keygray,colback=keybg,coltitle=white,colbacktitle=keygray,fonttitle=\bfseries,title={Question Recap},arc=3pt,breakable,boxrule=1pt}
\newtcolorbox{ideabox}{colframe=titleblue,colback=titleblue!6,coltitle=white,colbacktitle=titleblue,fonttitle=\bfseries,title={The Big Ideas},arc=3pt,breakable,boxrule=1.2pt}
\newtcolorbox{learnbox}{colframe=iagreen,colback=iagreen!5,coltitle=white,colbacktitle=iagreen,fonttitle=\bfseries,title={The Big Picture},arc=3pt,breakable,boxrule=1.2pt}
\newcommand{\lh}[1]{\par\smallskip{\bfseries\color{iagreen}#1}\par\nobreak\smallskip}

\newcommand{\soln}{\par\smallskip\textit{\textbf{Solution.}}\ }
\newcommand{\workspace}[1]{\par\vspace{3pt}\fcolorbox{black!30}{white}{\begin{minipage}[t][#1][t]{\dimexpr\linewidth-2\fboxsep\relax}\ \end{minipage}}\par}
\newcommand{\sech}[1]{\par\vspace{8pt}{\Large\bfseries\color{iagreen}#1}\par\nobreak\vspace{2pt}{\color{black!20}\hrule height1pt}\vspace{6pt}}

% A compact coordinate plot sized for an example box: \eplot{xmin}{xmax}{ymin}{ymax}{body}
\newcommand{\eplot}[5]{\begin{center}\begin{tikzpicture}
\begin{axis}[width=6.8cm,height=4.4cm,axis lines=middle,xlabel={\tiny$x$},ylabel={\tiny$y$},
grid=both,grid style={gray!16},xmin=#1,xmax=#2,ymin=#3,ymax=#4,
xtick distance=2,ytick distance=2,tick label style={font=\tiny},axis line style={-},samples=120]
#5
\end{axis}\end{tikzpicture}\end{center}}

% A sine-curve plot (degrees on x, 0..360): \splot{ymin}{ymax}{body}
\newcommand{\splot}[3]{\begin{center}\begin{tikzpicture}
\begin{axis}[width=8.6cm,height=4.6cm,axis lines=middle,xlabel={\tiny$x\,(^\circ)$},ylabel={\tiny$y$},
grid=both,grid style={gray!16},xmin=0,xmax=360,ymin=#1,ymax=#2,
xtick distance=90,ytick distance=1,tick label style={font=\tiny},axis line style={-},samples=180]
#3
\end{axis}\end{tikzpicture}\end{center}}

% A small coordinate plot: \plot{xmin}{xmax}{ymin}{ymax}{body}
\newcommand{\plot}[5]{\begin{center}\begin{tikzpicture}
\begin{axis}[width=8.4cm,height=6cm,axis lines=middle,xlabel={\scriptsize$x$},ylabel={\scriptsize$y$},
grid=both,grid style={gray!16},xmin=#1,xmax=#2,ymin=#3,ymax=#4,
xtick distance=2,ytick distance=2,tick label style={font=\tiny},axis line style={-}]
#5
\end{axis}\end{tikzpicture}\end{center}}

% Right triangle (angle theta at bottom-left, right angle at bottom-right):
%   \rtri{drawBase}{drawHeight}{oppLabel}{adjLabel}{hypLabel}
\newcommand{\rtri}[5]{\begin{center}\begin{tikzpicture}[scale=0.85]
\coordinate (A) at (0,0);\coordinate (B) at (#1,0);\coordinate (C) at (#1,#2);
\draw[thick,exblue] (A)--(B)--(C)--cycle;
\draw[black] ($(B)+(-0.32,0)$)--($(B)+(-0.32,0.32)$)--($(B)+(0,0.32)$);
\node[below] at ($(A)!0.5!(B)$){#4};
\node[right] at ($(B)!0.5!(C)$){#3};
\node[above left] at ($(A)!0.5!(C)$){#5};
\node at (0.7,0.22){$\theta$};
\end{tikzpicture}\end{center}}

% General (oblique) triangle with vertex labels and opposite-side labels:
%   \gtri{Alabel}{Blabel}{Clabel}{aSide}{bSide}{cSide}
\newcommand{\gtri}[6]{\begin{center}\begin{tikzpicture}[scale=0.85]
\coordinate (A) at (0,0);\coordinate (B) at (5,0);\coordinate (C) at (1.6,3);
\draw[thick,exblue] (A)--(B)--(C)--cycle;
\node[below left] at (A){#1};\node[below right] at (B){#2};\node[above] at (C){#3};
\node[below] at ($(A)!0.5!(B)$){#6};
\node[right] at ($(B)!0.5!(C)$){#4};
\node[left] at ($(A)!0.5!(C)$){#5};
\end{tikzpicture}\end{center}}

\fancyhead[R]{\footnotesize MCV4U \textperiodcentered{} 4.3}
\begin{document}

\begin{center}
{\LARGE\bfseries\color{titleblue}Worksheet 4.3 --- Critical, Inflection \& End Behaviour}\\[2pt]
{\small Grade 12 Calculus and Vectors (MCV4U) \textperiodcentered{} Unit 4: Curve Sketching}
\end{center}

\begin{ideabox}
Gather all features before sketching: domain, intercepts, asymptotes, extrema (from $f'$), inflection/concavity (from $f''$), and end behaviour.
\begin{itemize}[leftmargin=*,itemsep=2pt]
  \item From $f'$: increase/decrease and extrema.
  \item From $f''$: concavity and inflection.
  \item From $f$: domain, intercepts, asymptotes, end behaviour.
\end{itemize}
\end{ideabox}

\begin{learnbox}
\lh{Gathering the features}A full analysis collects critical numbers, inflection points, intercepts, and the behaviour as $x\to\pm\infty$.
\lh{Asymptotes and limits}Limits at infinity reveal horizontal asymptotes, while limits at breaks reveal vertical ones.
\lh{A complete picture}Together these features describe every important part of the graph before you draw it.
\end{learnbox}

\sech{Worked Examples}

\begin{example}{ex0}{Example 1 --- List features}
Key features of $f(x)=x^3-3x^2$.\soln
\textbf{Step 1 (find extrema from $f'$):} $f'(x)=3x^2-6x=3x(x-2)=0\Rightarrow x=0,2$; testing signs gives max $(0,0)$ and min $(2,-4)$.\\[3pt]
\textbf{Step 2 (find inflection from $f''$):} $f''(x)=6x-6=0\Rightarrow x=1$, giving $f(1)=-2$, so $(1,-2)$.\\[3pt]
\textbf{Step 3 (end behaviour):} odd-degree leading term $x^3$ with positive coefficient $\Rightarrow$ down on the left, up on the right.\\[3pt]
\textbf{Conclusion:} max $(0,0)$, min $(2,-4)$, inflection $(1,-2)$, ends down--up:\eplot{-1.5}{4}{-6}{4}{\addplot[exblue,very thick,domain=-1:3.3]{x^3-3*x^2};\addplot[red,only marks,mark=*,mark size=1.6pt] coordinates {(0,0) (2,-4) (1,-2)};}
\end{example}

\begin{example}{ex1}{Example 2 --- Read f from f'}
$f'>0$ on $(-\infty,1)$, $f'<0$ on $(1,\infty)$. Describe $f$ near $x=1$.\soln
\textbf{Step 1 (interpret the sign pattern):} $f'$ goes from positive to negative at $x=1$.\\[3pt]
\textbf{Step 2 (apply the first-derivative test):} a $+\to-$ change means $f$ switches from increasing to decreasing.\\[3pt]
\textbf{Conclusion:} local max at $x=1$.
\end{example}

\begin{example}{ex2}{Example 3 --- Classify from f''}
$f'(2)=0$ and $f''(2)=-3$. Classify $x=2$.\soln
\textbf{Step 1 (confirm it's a critical point):} $f'(2)=0$.\\[3pt]
\textbf{Step 2 (apply the second-derivative test):} $f''(2)=-3<0$, so the curve is concave down there.\\[3pt]
\textbf{Conclusion:} local maximum at $x=2$.
\end{example}

\begin{example}{ex3}{Example 4 --- Rational asymptotes}
Asymptotes of $f(x)=\dfrac{x}{x-2}$.\soln
\textbf{Step 1 (find the vertical asymptote, where the denominator is $0$):} $x-2=0\Rightarrow x=2$.\\[3pt]
\textbf{Step 2 (find the horizontal asymptote by comparing degrees):} numerator and denominator both have degree $1$, so the HA is the ratio of leading coefficients, $\tfrac11=1$.\\[3pt]
\textbf{Conclusion:} vertical asymptote $x=2$; horizontal asymptote $y=1$.
\end{example}

\begin{example}{ex4}{Example 5 --- Inflection, not extremum}
Inflection of $f(x)=x^3+1$; why no extremum?\soln
\textbf{Step 1 (find $f''$ and set to $0$):} $f''(x)=6x=0\Rightarrow x=0$, giving $f(0)=1$, so $(0,1)$.\\[3pt]
\textbf{Step 2 (check $f'$ to see if it's also an extremum):} $f'(x)=3x^2\ge0$ for every $x$ --- never negative, so no sign change.\\[3pt]
\textbf{Conclusion:} $(0,1)$ is an inflection point but not an extremum, since $f'$ never changes sign there.
\end{example}

\begin{example}{ex5}{Example 6 --- Critical numbers}
Critical numbers of $f(x)=x^3-3x^2$.\soln
\textbf{Step 1 (differentiate):} $f'(x)=3x^2-6x$.\\[3pt]
\textbf{Step 2 (factor and set to $0$):} $3x(x-2)=0$.\\[3pt]
\textbf{Conclusion:} $x=0$ and $x=2$.
\end{example}

\begin{example}{ex6}{Example 7 --- Classify from f''}
$f'(3)=0,\ f''(3)>0$. Classify $x=3$.\soln
\textbf{Step 1 (confirm it's a critical point):} $f'(3)=0$.\\[3pt]
\textbf{Step 2 (apply the second-derivative test):} $f''(3)>0$ means concave up there.\\[3pt]
\textbf{Conclusion:} local minimum at $x=3$.
\end{example}

\begin{example}{ex7}{Example 8 --- Inflection}
Inflection of $f(x)=x^3-3x^2$.\soln
\textbf{Step 1 (find $f''$):} $f''(x)=6x-6$.\\[3pt]
\textbf{Step 2 (set $f''=0$ and solve):} $x=1$, giving $f(1)=1-3=-2$.\\[3pt]
\textbf{Conclusion:} inflection point at $(1,-2)$.
\end{example}

\begin{example}{ex8}{Example 9 --- End behaviour}
End behaviour of $f(x)=-x^3+2x$.\soln
\textbf{Step 1 (identify the leading term):} the highest-degree term is $-x^3$, degree $3$ (odd) with negative coefficient.\\[3pt]
\textbf{Step 2 (recall the pattern for odd-degree, negative leading coefficient):} the graph rises on the far left and falls on the far right.\\[3pt]
\textbf{Conclusion:} up--down end behaviour.
\end{example}

\clearpage
\sech{Your Turn --- Show Your Work}
For each question, show all steps clearly in the space provided.

\begin{qbox}{Question 1}
Critical numbers of $f(x)=x^3-3x^2$?\workspace{2.4cm}
\end{qbox}

\begin{qbox}{Question 2}
If $f'(3)=0$ and $f''(3)>0$, classify $x=3$.\workspace{2.4cm}
\end{qbox}

\begin{qbox}{Question 3}
Inflection point of $f(x)=x^3-3x^2$?\workspace{2.4cm}
\end{qbox}

\begin{qbox}{Question 4}
Asymptotes of $f(x)=\dfrac1{x-3}$?\workspace{2.4cm}
\end{qbox}

\begin{qbox}{Question 5}
End behaviour of $f(x)=-x^3+2x$?\workspace{2.4cm}
\end{qbox}

\begin{qbox}{Question 6}
Local extrema of $f(x)=x^3-3x$?\workspace{2.4cm}
\end{qbox}

\begin{qbox}{Question 7}
Asymptotes of $f(x)=\dfrac{x}{x-2}$?\workspace{2.4cm}
\end{qbox}

\begin{qbox}{Question 8}
If $f'(1)=0$ and $f''(1)<0$, classify $x=1$.\workspace{2.4cm}
\end{qbox}

\begin{qbox}{Question 9}
Inflection point of $f(x)=x^3$?\workspace{2.4cm}
\end{qbox}

\begin{qbox}{Question 10}
End behaviour of $f(x)=x^4-1$?\workspace{2.4cm}
\end{qbox}

\begin{qbox}{Question 11}
Where is $f(x)=x^3-3x^2$ increasing?\workspace{2.4cm}
\end{qbox}

\begin{qbox}{Question 12}
Does $f(x)=x^3+1$ have a local extremum?\workspace{2.4cm}
\end{qbox}

\begin{qbox}{Question 13 --- Challenge}
List every key feature of $f(x)=x^3-3x$ (intercepts, extrema, inflection, ends).\workspace{4cm}
\end{qbox}

\clearpage
\sech{Question Recap \& Answer Key}

\begin{recapbox}
\begin{multicols}{2}
\small
\begin{enumerate}[leftmargin=*,itemsep=3pt]
  \item Critical numbers of $f(x)=x^3-3x^2$?
  \item If $f'(3)=0$ and $f''(3)>0$, classify $x=3$.
  \item Inflection point of $f(x)=x^3-3x^2$?
  \item Asymptotes of $f(x)=\dfrac1{x-3}$?
  \item End behaviour of $f(x)=-x^3+2x$?
  \item Local extrema of $f(x)=x^3-3x$?
  \item Asymptotes of $f(x)=\dfrac{x}{x-2}$?
  \item If $f'(1)=0$ and $f''(1)<0$, classify $x=1$.
  \item Inflection point of $f(x)=x^3$?
  \item End behaviour of $f(x)=x^4-1$?
  \item Where is $f(x)=x^3-3x^2$ increasing?
  \item Does $f(x)=x^3+1$ have a local extremum?
  \item List every key feature of $f(x)=x^3-3x$ (intercepts, extrema, inflection, ends).
\end{enumerate}
\end{multicols}
\end{recapbox}

\begin{keybox}
\begin{multicols}{2}
\small
\begin{enumerate}[leftmargin=*,itemsep=3pt]
  \item $x=0,2$
  \item local min
  \item $(1,-2)$
  \item VA $x=3$, HA $y=0$
  \item up–down
  \item max $(-1,2)$, min $(1,-2)$
  \item VA $x=2$, HA $y=1$
  \item local max
  \item $(0,0)$
  \item up–up
  \item $x<0$ or $x>2$
  \item no (inflection at $(0,1)$)
  \item x-int $0,\pm\sqrt3$; max $(-1,2)$; min $(1,-2)$; inflection $(0,0)$; ends down–up
\end{enumerate}
\end{multicols}
\end{keybox}

\end{document}
